% modul-2.tex - Modul 2: Python-Grundlagen mit KI verstehen

\startcomponent modul-2
\environment ../pyt-layout
\environment ../pyt-env

\chapter{Python-Grundlagen mit KI verstehen}

{\it Dauer: 1 Tag (4 Lektionen à 50 Minuten) • Voraussetzung: Modul 1 abgeschlossen}

\startlernziele
Nach diesem Modul können Sie:
\startitemize[n,packed]
    \item Grundlegende Python-Datentypen verstehen und anwenden
    \item Kontrollstrukturen (if/else, Schleifen) implementieren
    \item Eigene Funktionen mit Parametern erstellen
    \item Mit Listen und Dictionaries arbeiten
    \item Python-Module importieren und nutzen
    \item KI-generierten Code evaluieren und anpassen
    \item Systematisch Fehler finden und beheben
\stopitemize
\stoplernziele

% =============================================================================
% VORBEREITUNG
% =============================================================================

\section{Vorbereitung}

{\it Zeitaufwand: 2-3 Stunden vor dem Präsenzunterricht}

\subsection{Leseauftrag: Python-Syntax Grundlagen}

\subsubsection{Variablen und Datentypen}

Lesen Sie folgende Konzepte und lassen Sie sich von der KI Beispiele erklären:

\startitemize[packed]
    \item {\bf int} -- Ganze Zahlen: \code{alter = 25}
    \item {\bf float} -- Dezimalzahlen: \code{preis = 19.99}
    \item {\bf str} -- Text/Strings: \code{name = "Anna"}
    \item {\bf bool} -- Wahrheitswerte: \code{ist_aktiv = True}
    \item {\bf None} -- Kein Wert: \code{ergebnis = None}
\stopitemize

\subsubsection{Operatoren}

\starttabulate[|l|l|l|]
\HL
\NC {\bf Typ} \NC {\bf Operatoren} \NC {\bf Beispiel} \NC\NR
\HL
\NC Arithmetisch \NC +, -, *, /, //, \%, ** \NC \code{5 + 3 = 8} \NC\NR
\NC Vergleich \NC ==, !=, <, >, <=, >= \NC \code{5 > 3} → True \NC\NR
\NC Logisch \NC and, or, not \NC \code{True and False} \NC\NR
\HL
\stoptabulate

\subsection{REPL-Experimente}

Öffnen Sie ein Terminal und starten Sie Python:

\startbashcode
python
\stopbashcode

Experimentieren Sie mit folgenden Befehlen:

\startpythoncode
# Variablen erstellen
name = "Anna"
alter = 25

# Operationen
print(f"Hallo {name}, du bist {alter} Jahre alt")
print(f"In 10 Jahren bist du {alter + 10}")

# Listen
zahlen = [1, 2, 3, 4, 5]
print(sum(zahlen))  # 15
print(len(zahlen))  # 5

# Dictionary
person = {"name": "Anna", "alter": 25}
print(person["name"])  # Anna
\stoppythoncode

\subsection{KI-Konzepterklärungen}

Stellen Sie der KI folgende Fragen:

\startitemize[n,packed]
    \item "Erkläre mir den Unterschied zwischen \code{=} und \code{==} in Python"
    \item "Was ist der Unterschied zwischen einer Liste und einem Dictionary?"
    \item "Wann verwende ich \code{for} und wann \code{while}?"
    \item "Was bedeutet \code{def} in Python?"
\stopitemize

% =============================================================================
% PRAXIS
% =============================================================================

\section{Praxis}

{\it Präsenzunterricht: 4 Lektionen à 50 Minuten}

\subsection{Lektion 1: Variablen, Datentypen \& Operatoren}

{\it Dauer: 50 Minuten}

\subsubsection{Variablen}

Variablen sind "Behälter" für Werte:

\startpythoncode
# Variable erstellen
name = "Anna"
alter = 25
groesse = 1.75
ist_student = True
\stoppythoncode

{\bf Naming Conventions (PEP 8):}

\startpythoncode
# Gut - snake_case
benutzer_name = "Anna"
max_versuche = 3

# Schlecht
BenutzerName = "Anna"  # CamelCase nur fuer Klassen
2name = "Bob"          # Darf nicht mit Zahl beginnen
\stoppythoncode

\subsubsection{Datentypen}

\startpythoncode
# Integer (Ganzzahlen)
alter = 25
anzahl = 100

# Float (Dezimalzahlen)
preis = 19.99
pi = 3.14159

# String (Text)
name = "Anna"
text = f"Ich bin {alter} Jahre alt"  # F-String

# Boolean (Wahrheitswert)
ist_aktiv = True
ist_fertig = False
\stoppythoncode

\subsubsection{Type Hints}

\startpythoncode
# Type Hints fuer bessere Lesbarkeit
name: str = "Anna"
alter: int = 25
groesse: float = 1.75
ist_student: bool = True

def begruessung(name: str) -> str:
    return f"Hallo {name}"
\stoppythoncode

\subsubsection{Type Conversion}

\startpythoncode
# String zu Zahl
alter_str = "25"
alter_int = int(alter_str)      # 25
preis_float = float("19.99")    # 19.99

# Zahl zu String
zahl = 42
text = str(zahl)  # "42"

# Fehlerbehandlung
try:
    zahl = int("abc")  # ValueError!
except ValueError:
    print("Keine gueltige Zahl")
\stoppythoncode


\subsection{Lektion 2: Kontrollstrukturen}

{\it Dauer: 50 Minuten}

\subsubsection{if/elif/else}

\startpythoncode
alter = 18

if alter >= 18:
    print("Volljaehrig")
else:
    print("Minderjaehrig")

# Mehrere Bedingungen
punkte = 85

if punkte >= 90:
    note = "A"
elif punkte >= 80:
    note = "B"
elif punkte >= 70:
    note = "C"
else:
    note = "F"
\stoppythoncode

\subsubsection{for-Schleifen}

\startpythoncode
# Ueber Liste iterieren
namen = ["Anna", "Bob", "Clara"]

for name in namen:
    print(f"Hallo {name}")

# Mit range()
for i in range(5):      # 0, 1, 2, 3, 4
    print(i)

for i in range(1, 6):   # 1, 2, 3, 4, 5
    print(i)

# Mit enumerate()
for index, name in enumerate(namen):
    print(f"{index + 1}. {name}")
\stoppythoncode

\subsubsection{while-Schleifen}

\startpythoncode
zaehler = 0

while zaehler < 5:
    print(zaehler)
    zaehler += 1

# Mit break
while True:
    antwort = input("Weiter? (ja/nein): ")
    if antwort == "nein":
        break
\stoppythoncode

\subsubsection{break, continue, pass}

\startpythoncode
# break - Schleife beenden
for i in range(10):
    if i == 5:
        break
    print(i)  # 0, 1, 2, 3, 4

# continue - Iteration ueberspringen
for i in range(10):
    if i % 2 == 0:
        continue
    print(i)  # 1, 3, 5, 7, 9
\stoppythoncode

\subsection{Lektion 3: Listen \& Dictionaries}

{\it Dauer: 50 Minuten}

\subsubsection{Listen}

\startpythoncode
# Liste erstellen
zahlen = [1, 2, 3, 4, 5]
namen = ["Anna", "Bob", "Clara"]
gemischt = [1, "zwei", 3.0, True]

# Zugriff (0-basierter Index)
print(namen[0])   # Anna
print(namen[-1])  # Clara (letztes Element)

# Liste aendern
namen.append("David")      # Hinzufuegen
namen.remove("Bob")        # Entfernen
namen.insert(0, "Zoe")     # An Position einfuegen

# Slicing
zahlen = [0, 1, 2, 3, 4, 5]
print(zahlen[1:4])    # [1, 2, 3]
print(zahlen[:3])     # [0, 1, 2]
print(zahlen[3:])     # [3, 4, 5]
\stoppythoncode

\subsubsection{List Comprehensions}

\startpythoncode
# Traditionell
quadrate = []
for x in range(5):
    quadrate.append(x ** 2)

# Mit List Comprehension
quadrate = [x ** 2 for x in range(5)]  # [0, 1, 4, 9, 16]

# Mit Bedingung
gerade = [x for x in range(10) if x % 2 == 0]  # [0, 2, 4, 6, 8]
\stoppythoncode

\subsubsection{Dictionaries}

\startpythoncode
# Dictionary erstellen
person = {
    "name": "Anna",
    "alter": 25,
    "stadt": "Zuerich"
}

# Zugriff
print(person["name"])           # Anna
print(person.get("alter", 0))   # 25 (mit Default)

# Aendern/Hinzufuegen
person["email"] = "anna@example.com"
person["alter"] = 26

# Iterieren
for key, value in person.items():
    print(f"{key}: {value}")
\stoppythoncode

\subsection{Lektion 4: Funktionen \& Module}

{\it Dauer: 50 Minuten}

\subsubsection{Funktionen definieren}

\startpythoncode
def begruessung(name: str) -> str:
    """Gibt eine Begruessung zurueck."""
    return f"Hallo {name}!"

# Aufruf
nachricht = begruessung("Anna")
print(nachricht)  # Hallo Anna!
\stoppythoncode

\subsubsection{Parameter und Rückgabewerte}

\startpythoncode
# Default-Parameter
def begruessung(name: str, sprache: str = "de") -> str:
    if sprache == "de":
        return f"Hallo {name}!"
    else:
        return f"Hello {name}!"

# Mehrere Rueckgabewerte
def statistik(zahlen: list) -> tuple:
    return min(zahlen), max(zahlen), sum(zahlen) / len(zahlen)

minimum, maximum, durchschnitt = statistik([1, 2, 3, 4, 5])
\stoppythoncode

\subsubsection{Module importieren}

\startpythoncode
# Standard-Module
import math
print(math.sqrt(16))  # 4.0
print(math.pi)        # 3.14159...

import random
print(random.randint(1, 10))  # Zufallszahl 1-10

from datetime import datetime
print(datetime.now())  # Aktuelles Datum/Zeit

# Eigene Module
# Datei: mathe_utils.py
# from mathe_utils import addiere
\stoppythoncode



% =============================================================================
% ÜBUNGEN
% =============================================================================

\section{Übungen}

\startaufgabenbox
{\bf Übung 1: Altersrechner / Währungsrechner} {\it (15 Min.)}

{\bf Option A: Altersrechner}
\startitemize[packed]
    \item Geburtsjahr einlesen
    \item Alter berechnen (aktuelles Jahr: 2025)
    \item Alter in Tagen berechnen (365 Tage/Jahr)
\stopitemize

{\bf Option B: Währungsrechner}
\startitemize[packed]
    \item Betrag in EUR eingeben
    \item Umrechnung in CHF (Kurs: 1 EUR = 0.95 CHF)
    \item Ergebnis auf 2 Dezimalstellen
\stopitemize
\stopaufgabenbox

\blank[medium]

\startaufgabenbox
{\bf Übung 2: Primzahl-Checker / Menü-System} {\it (15 Min.)}

{\bf Option A: Primzahl-Checker}
\startitemize[packed]
    \item Zahl vom Benutzer einlesen
    \item Prüfen, ob Primzahl
    \item Programm läuft bis 'quit'
\stopitemize

{\bf Option B: Menü-System}
\startitemize[packed]
    \item Menü mit 3 Optionen anzeigen
    \item Benutzerwahl einlesen
    \item Läuft bis "Beenden" gewählt wird
\stopitemize
\stopaufgabenbox

\blank[medium]

\startaufgabenbox
{\bf Übung 3: Inventar-System / Duplikat-Entferner} {\it (15 Min.)}

{\bf Option A: Inventar-System}
\startitemize[packed]
    \item Dictionary für Artikel + Anzahl
    \item Artikel hinzufügen/entfernen
    \item Gesamtbestand anzeigen
\stopitemize

{\bf Option B: Duplikat-Entferner}
\startitemize[packed]
    \item Liste mit Duplikaten erstellen
    \item Duplikate entfernen (Reihenfolge beibehalten)
    \item Beide Listen ausgeben
\stopitemize
\stopaufgabenbox

\blank[medium]

\startaufgabenbox
{\bf Übung 4: Utility-Funktionen-Modul} {\it (20 Min.)}

Erstellen Sie ein Modul \code{utils.py} mit folgenden Funktionen:

\startpythoncode
def ist_palindrom(text: str) -> bool:
    """Prueft, ob Text ein Palindrom ist."""
    ...

def wort_zaehlen(text: str) -> int:
    """Zaehlt die Woerter im Text."""
    ...

def durchschnitt(zahlen: list) -> float:
    """Berechnet den Durchschnitt."""
    ...
\stoppythoncode
\stopaufgabenbox

% =============================================================================
% NACHBEARBEITUNG
% =============================================================================

\section{Nachbearbeitung}

{\it Hausaufgaben: 6-7 Stunden vor Modul 3}

\subsection{Aufgabe 1: Persönlicher Finanztracker}

Erstellen Sie einen einfachen Finanztracker:

\startitemize[packed]
    \item Einnahmen und Ausgaben erfassen
    \item Kategorien (Essen, Transport, Freizeit, ...)
    \item Bilanz berechnen
    \item Daten in Dictionary speichern
    \item Ausgabe formatieren
\stopitemize

\startpythoncode
# Beispiel-Struktur
finanzen = {
    "einnahmen": [
        {"betrag": 3000, "kategorie": "Gehalt"},
        {"betrag": 100, "kategorie": "Verkauf"}
    ],
    "ausgaben": [
        {"betrag": 50, "kategorie": "Essen"},
        {"betrag": 200, "kategorie": "Miete"}
    ]
}
\stoppythoncode

\subsection{Aufgabe 2: Code-Comprehension}

Analysieren Sie folgenden Code und erklären Sie ihn Zeile für Zeile:

\startpythoncode
def fibonacci(n: int) -> list:
    if n <= 0:
        return []
    elif n == 1:
        return [0]

    fib = [0, 1]
    while len(fib) < n:
        fib.append(fib[-1] + fib[-2])
    return fib

ergebnis = fibonacci(10)
gerade = [x for x in ergebnis if x % 2 == 0]
summe = sum(gerade)
\stoppythoncode

\subsection{Aufgabe 3: Algorithm Implementation}

Implementieren Sie folgende Algorithmen:

\startitemize[n,packed]
    \item {\bf Bubble Sort} -- Liste aufsteigend sortieren
    \item {\bf Binary Search} -- Element in sortierter Liste finden
    \item {\bf Fakultät} -- n! berechnen (iterativ und rekursiv)
\stopitemize

\subsection{Aufgabe 4: Real-World Mini-Projekt}

Wählen Sie eines:

\startitemize[packed]
    \item {\bf Todo-Listen-Manager} -- Aufgaben verwalten mit Prioritäten
    \item {\bf Quiz-App} -- Fragen stellen, Antworten prüfen, Punkte zählen
    \item {\bf Kontaktbuch} -- Kontakte speichern, suchen, bearbeiten
\stopitemize

\subsection{Reflexion}

Beantworten Sie schriftlich (400+ Wörter):
\startitemize[n,packed]
    \item Welches Konzept war am schwierigsten zu verstehen?
    \item Wie hat KI beim Lernen geholfen?
    \item Was würden Sie beim nächsten Mal anders machen?
    \item Welche Programmierprojekte möchten Sie umsetzen?
\stopitemize

% =============================================================================
% MATERIALIEN
% =============================================================================

\section{Materialien}

\subsection{Python-Syntax Quick Reference}

\starttabulate[|l|p(9cm)|]
\HL
\NC {\bf Element} \NC {\bf Syntax} \NC\NR
\HL
\NC Variable \NC \code{name = "Wert"} \NC\NR
\NC F-String \NC \code{f"Text \{variable\}"} \NC\NR
\NC Liste \NC \code{[1, 2, 3]} \NC\NR
\NC Dictionary \NC \code{\{"key": "value"\}} \NC\NR
\NC Funktion \NC \code{def name(param): return wert} \NC\NR
\NC If \NC \code{if bedingung: ...} \NC\NR
\NC For \NC \code{for x in liste: ...} \NC\NR
\NC While \NC \code{while bedingung: ...} \NC\NR
\HL
\stoptabulate

\subsection{Debugging-Guide}

{\bf Häufige Fehler und Lösungen:}

\starttabulate[|l|p(8cm)|]
\HL
\NC {\bf Fehler} \NC {\bf Lösung} \NC\NR
\HL
\NC SyntaxError \NC Klammern, Doppelpunkte, Anführungszeichen prüfen \NC\NR
\NC NameError \NC Variable definiert? Schreibweise korrekt? \NC\NR
\NC TypeError \NC Datentypen kompatibel? Conversion nötig? \NC\NR
\NC IndexError \NC Index innerhalb der Liste? len() prüfen \NC\NR
\NC KeyError \NC Schlüssel existiert? .get() verwenden \NC\NR
\HL
\stoptabulate

{\bf Debugging-Strategien:}
\startitemize[n,packed]
    \item \code{print()} an kritischen Stellen
    \item Fehlermeldung genau lesen
    \item Code schrittweise testen
    \item KI nach Erklärung fragen
\stopitemize

\blank[big]

\starttippbox
{\bf Erfolgskriterien für Modul 2}

Sie haben das Modul erfolgreich abgeschlossen, wenn Sie:
\startitemize[packed]
    \item Alle grundlegenden Python-Datentypen verstehen
    \item Kontrollstrukturen implementieren können
    \item Funktionen mit Parametern erstellen können
    \item Mit Listen und Dictionaries arbeiten können
    \item Einfache Programme debuggen können
    \item Alle Nachbearbeitungsaufgaben eingereicht haben
\stopitemize
\stoptippbox

\stopcomponent